\chapter{统计与概率}\label{ch:g9:statproba}

数据与机会无处不在：体育比赛结果、天气预报，
游戏。本章介绍如何用平均数、中位数和极差概括
一组数，以及如何计算简单的
随机试验的概率 --- 包括两步试验，
借助树状图处理。两个主题在本系列的
高中分册中会有更深入的发展。

\section{数据的概括}

\begin{definition}[平均数、中位数、极差]\label{def:g9:statproba:indicators}
对于 $N$ 个数构成的一组数：
\begin{itemize}
  \item \emph{平均数}\index{平均数}是所有数值之和
        除以 $N$；
  \item \emph{中位数}\index{中位数}是将排序后的一组数分成个数相等的两半的数值：
        当 $N$ 为奇数时，为中间的数值；
        当 $N$ 为偶数时，为中间两个数值的中点；
  \item \emph{极差}\index{极差}是最大值
        减去最小值。
\end{itemize}
\end{definition}

\begin{example}\label{ex:g9:statproba:indicators}
某学生的成绩为 $8,\ 12,\ 9,\ 15,\ 11$。
\begin{enumerate}
  \item 平均数：$\dfrac{8 + 12 + 9 + 15 + 11}{5} = \dfrac{55}{5} = 11$。
  \item 中位数：将这组数排序得 $8, 9, 11, 12, 15$；中间（第 3 个）
        数值为 $11$。
  \item 极差：$15 - 8 = 7$。
\end{enumerate}
若再增加一个成绩 $17$，排序后为 $8, 9, 11, 12, 15, 17$，中位数
变为 $11$ 与 $12$ 的中点，即 $11.5$，平均数
变为 $\frac{72}{6} = 12$。
\end{example}

\begin{example}[加权平均数]\label{ex:g9:statproba:weighted}
某日售出的鞋码如下：
\begin{center}
\renewcommand{\arraystretch}{1.2}
\begin{tabular}{c|ccccc}
鞋码 & $38$ & $39$ & $40$ & $41$ & $42$ \\
\hline
频数 & $3$ & $7$ & $6$ & $3$ & $1$ \\
\end{tabular}
\end{center}
平均鞋码为
\[
\frac{3 \times 38 + 7 \times 39 + 6 \times 40 + 3 \times 41 + 1 \times 42}
{3 + 7 + 6 + 3 + 1}
= \frac{114 + 273 + 240 + 123 + 42}{20}
= \frac{792}{20} = 39.6 .
\]
\end{example}

\begin{remark}
平均数与中位数可能相差很大。在数列
$10, 10, 10, 10, 60$ 中，平均数为 $20$（被数值 $60$ 拉高），
而中位数为 $10$。应当总是思考哪一个指标
更能代表实际情况。
\end{remark}

\section{概率}

\begin{definition}[概率]\label{def:g9:statproba:probability}
一个随机试验可能有若干种\emph{结果}；
\emph{事件}是若干结果构成的集合。每个结果都有一个
\emph{概率}\index{概率}：它是介于 $0$ 与 $1$ 之间的数，
所有结果的概率之和为 $1$；
$\P(A)$ 为事件 $A$ 的概率，即其所含各结果概率之和。
当全部 $n$ 个结果等可能时，
\[
\P(A) = \frac{\text{有利结果数}}{n}.
\]
\end{definition}

\begin{example}\label{ex:g9:statproba:spinner}
一个转盘被等分为 $8$ 个扇形：$3$ 个红色、$3$ 个蓝色、$2$ 个绿色。
每个扇形的概率为 $\frac18$，因此
\[
\P(\text{红色}) = \frac38, \qquad
\P(\text{绿色}) = \frac28 = \frac14, \qquad
\P(\text{非绿色}) = 1 - \frac14 = \frac34 .
\]
\end{example}

\begin{proposition}[对立事件]\label{prop:g9:statproba:complement}
对于任一事件 $A$，其\emph{对立事件} $\bar A$（“$A$
不发生”）满足
\[
\P(\bar A) = 1 - \P(A).
\]
\end{proposition}

\begin{proof}
每个结果恰属于 $A$ 与 $\bar A$ 之一，因此
$\P(A) + \P(\bar A)$ 即所有结果的概率之和，其值为
$1$。
\end{proof}

\section{两步试验}

\begin{method}[树状图]\label{met:g9:statproba:tree}
对于两步试验：
\begin{enumerate}
  \item 对第一步的每个结果画一条分支，然后在各分支之后接上
        第二步的各个结果，并
        在每条分支上标注其概率；
  \item 沿一条路径将各概率相乘，
        得到该路径的概率；
  \item 将构成某一事件的所有路径的概率相加。
\end{enumerate}
\end{method}

\begin{example}\label{ex:g9:statproba:tree}
袋中装有 $2$ 个红色、$3$ 个黑色的小球。取出一个小球，记下它的
颜色，\emph{放回}，再取一次。每次取球，得红色的概率为 $\frac25$，
得黑色的概率为 $\frac35$。

\begin{omfigure}
\begin{tikzpicture}[scale=0.95]
  \coordinate (root) at (0,0);
  \coordinate (R) at (2.2,1.0);
  \coordinate (B) at (2.2,-1.0);
  \coordinate (RR) at (4.7,1.55);
  \coordinate (RB) at (4.7,0.45);
  \coordinate (BR) at (4.7,-0.45);
  \coordinate (BB) at (4.7,-1.55);
  \draw (root) -- (R) node[midway, above, font=\small, sloped] {$\frac25$};
  \draw (root) -- (B) node[midway, below, font=\small, sloped] {$\frac35$};
  \draw (R) -- (RR) node[midway, above, font=\small, sloped] {$\frac25$};
  \draw (R) -- (RB) node[midway, below, font=\small, sloped] {$\frac35$};
  \draw (B) -- (BR) node[midway, above, font=\small, sloped] {$\frac25$};
  \draw (B) -- (BB) node[midway, below, font=\small, sloped] {$\frac35$};
  \node[omThm, font=\small, above left=-2pt] at (R) {红};
  \node[font=\small, below left=-2pt] at (B) {黑};
  \node[right, font=\small] at (RR)
    {红红：\ $\frac25 \times \frac25 = \frac{4}{25}$};
  \node[right, font=\small] at (RB)
    {红黑：\ $\frac25 \times \frac35 = \frac{6}{25}$};
  \node[right, font=\small] at (BR)
    {黑红：\ $\frac35 \times \frac25 = \frac{6}{25}$};
  \node[right, font=\small] at (BB)
    {黑黑：\ $\frac35 \times \frac35 = \frac{9}{25}$};
\end{tikzpicture}

{\small 有放回的两次取球的树状图：沿各分支将概率相乘，
并检验四条路径的概率之和为 $1$。}
\end{omfigure}

取出两个颜色相同的小球的概率：
$\frac{4}{25} + \frac{9}{25} = \frac{13}{25}$。至少取出一个红球的概率为
$1 - \P(\text{黑黑}) = 1 - \frac{9}{25} = \frac{16}{25}$。
\end{example}

\begin{remark}[频率趋近于概率]
掷一枚均匀骰子 $6000$ 次，得到六点的次数\emph{大约}为 $1000$，而
并非恰好如此。随着重复试验次数增多，观测到的频率
越来越接近概率：正是这一点使概率成为
描述重复试验的恰当工具（这一内容在本系列的
高中分册中展开）。
\end{remark}

\section{练习}

\begin{exercise}[$\star$]\label{exo:g9:statproba:1}
求数据
$14,\ 9,\ 17,\ 9,\ 12,\ 11,\ 12$ 的平均数、中位数与极差。
\end{exercise}

\begin{exercise}[$\star$]\label{exo:g9:statproba:2}
一周中午的气温（单位：度）依次为 $19,\ 21,\ 24,\ 18,\ 22,\ 25,\
19$。求平均数（精确到十分位）和中位数。
\end{exercise}

\begin{exercise}[$\star$]\label{exo:g9:statproba:3}
某球队在 $10$ 场比赛中的进球数如下：
\begin{center}
\renewcommand{\arraystretch}{1.2}
\begin{tabular}{c|cccc}
进球数 & $0$ & $1$ & $2$ & $3$ \\
\hline
场数 & $2$ & $4$ & $3$ & $1$ \\
\end{tabular}
\end{center}
求每场比赛的平均进球数与中位数。
\end{exercise}

\begin{exercise}[$\star$]\label{exo:g9:statproba:4}
掷一枚均匀骰子一次，求下列事件的概率：“掷出
$3$ 点”；“掷出偶数”；“掷出至少 $3$ 点”；“不掷出
$6$ 点”。
\end{exercise}

\begin{exercise}[$\star$]\label{exo:g9:statproba:5}
盒中装有 $12$ 个球：白色 $5$ 个、红色 $4$ 个、绿色 $3$ 个。从中
任取一个球。求 $\P(\text{白色})$、$\P(\text{红色或绿色})$ 与
$\P(\text{非白色})$。后两个概率之间有什么关系？
\end{exercise}

\begin{exercise}[$\star\star$]\label{exo:g9:statproba:6}
掷一枚硬币，再掷一枚均匀骰子。画出树状图（或列举
所有结果），求出现正面\emph{且}掷出
$6$ 点的概率；再求出现正面\emph{或}掷出 $6$ 点（或两者都出现）的概率。
\end{exercise}

\begin{exercise}[$\star\star$]\label{exo:g9:statproba:7}
在\cref{ex:g9:statproba:tree}所述的袋中（内有 $2$ 个红色、$3$ 个黑色小球），两次
取球改为\emph{不放回}地进行。画出新的树状图
（注意：第二步的概率会改变），并求
两个都是黑色小球的概率，再求两个小球
颜色不同的概率。
\end{exercise}

\begin{exercise}[$\star\star$]\label{exo:g9:statproba:8}
某篮球运动员在 $9$ 场比赛后，场均得 $14$ 分。
\begin{enumerate}
  \item 她共得了多少分？
  \item 她在第 $10$ 场比赛中需得多少分，才能使场均得分达到
        $15$ 分？
\end{enumerate}
\end{exercise}

\begin{exercise}[$\star\star\star$]\label{exo:g9:statproba:9}
游戏规则：同时掷两枚均匀骰子，若两次结果相同则
获胜（即“掷出一对相同点数”）。
\begin{enumerate}
  \item 求获胜的概率。
  \item 游戏进行两次（两次独立），利用树状图求
        至少获胜一次的概率。
\end{enumerate}
\end{exercise}

\section{问题：梅雷骑士的破产赌注}

\begin{problem}[{周末问题 --- 1654 年开创概率论的
赌博争议，以及
骗过每间教室的生日巧合}]
\label{pb:g9:statproba:1}
1654 年，一位嗜赌的法国贵族梅雷骑士，
靠一次掷骰子的赌注发了财，又在另一次他认为等价的赌注上开始赔钱。
他深感困惑，便向数学家
布莱兹·帕斯卡请教，帕斯卡又写信给皮埃尔·德·费马，两人
的通信奠定了概率论的基础。本题借助本章的工具，
重演整个事件
——等可能结果、树状图与补事件法则
（\cref{prop:g9:statproba:complement}），并以
所有巧合中最违反直觉的一个作结。

\medskip
\noindent\textbf{第一部分 --- 诚实地数骰子。}
\begin{enumerate}
  \item 掷两次均匀骰子。用树状图（或 $6 \times 6$
        表格），求两次掷出中\emph{没有}六点的概率，
        并由此求出至少掷出一次
        六点的概率。
  \item 一位同学争辩道：“每掷一次有六分之一的机会，掷两次，
        所以 $\frac16 + \frac16 = \frac13$。”请与第
        1 题比较，指出他的加法把哪些确切结果
        计算了两次。
  \item 掷两枚骰子并求点数之和。求和为 $7$ 与和为 $2$ 的概率，
        并说明为什么 $7$ 最受
        赌徒青睐？
  \item 旧有争议：掷两枚骰子时，点数和为 $9$ 与点数和为
        $10$ 哪一个更可能？请数出所有结果，判定此争议。
  \item 推广第 1 题：说明为什么 $n$ 次掷骰中至少
        出现一次六点的概率为
        \[
        1 - \left(\frac56\right)^{\!n} ,
        \]
        树状图中“不出现六点”的分支逐层相乘，
        由此可得上式。
\end{enumerate}

\medskip
\noindent\textbf{第二部分 --- 梅雷骑士的两次赌注。}
\begin{enumerate}[resume]
  \item 第一次赌注（平赔）：\emph{掷一枚骰子四次，至少出现
        一次六点}。计算
        $\left(\frac56\right)^4$ 的分数形式与小数形式，
        以及梅雷骑士的获胜概率，
        并判断这一赌注是否划算？
  \item 梅雷骑士自己的推理是：“四次掷骰，每次为
        $\frac16$，共有 $\frac46$ 的机会。”这一推理
        在此得出了几乎正确的答案，但若推广到 $7$ 次掷骰，
        会导出什么荒谬的结论？它重复了第 2
        题中的哪一个错误？
  \item 第二次赌注，按比例推理认为与前者等价
        （“掷 $24$ 次，每次 $\frac1{36}$，同样得出
        $\frac{24}{36} = \frac46$”）：\emph{掷两枚骰子
        $24$ 次，至少出现一次双六}。计算真实的
        获胜概率 $1 - \left(\frac{35}{36}\right)^{24}$
        （用计算器）。梅雷骑士为何慢慢破产？
  \item 掷两枚骰子多少次才能使他恢复优势？
        用计算器检验 $n = 25$ 的情形，
        并作出结论。
  \item 用两句话概括其中的教训：一般而言，把概率乘以
        尝试次数，错在哪里？
        ——哪一个正确的工具能够取代这种
        诱惑？
\end{enumerate}

\medskip
\noindent\textbf{第三部分 --- 平均数、中位数与生日。}
\begin{enumerate}[resume]
  \item 一家小公司每月付给九名员工 $2\,000$ 欧元，
        付给经理 $20\,000$ 欧元。求工资的平均数与中位数
        （见\cref{def:g9:statproba:indicators}）。
        招聘广告应当如实引用
        哪一个数字？
  \item 构造一组五个考试成绩，使其平均数为 $12$，
        中位数为 $15$。“平均数小于中位数”说明
        成绩分布呈怎样的形状？
  \item 生日问题：在一个有 $23$ 名学生的班级中，
        至少有两人生日相同的概率是多少？
        列出补事件（$23$ 人的生日两两不同）：
        \[
        \frac{365}{365} \times \frac{364}{365} \times
        \frac{363}{365} \times \dots \times \frac{343}{365} .
        \]
        计算前三个因式的乘积，描述各因式的变化方式，
        并结合全部乘积约为 $0.493$ 这一事实，回答
        这一问题；多数人期望“非常不可能”，
        数学告诉我们的
        结论是什么？
  \item 修正直觉：在 $23$ 人的班级中，有多少\emph{对}学生
        可能生日相同？
        （即\cref{pb:g6:wholes:1}中的握手计数）？用一句话
        解释为什么真正制造惊讶的是这个数，而不是
        $23$ 本身。
  \item 结语：设计一个实验来检验任一结果
        ——用真实骰子重做梅雷骑士的赌注，或在若干班级中
        检验生日问题；并谨慎地说明，
        随着重复次数不断增加，你对观测频率
        有何期望？（该期望有名称，称为
        大数定律；其证明留待高中
        分册中给出。）
\end{enumerate}
\end{problem}
